%% ============================================================================
%%  ejemplo_capas.tex · Póster de ejemplo con la plantilla postercientifico.sty v1.0.0
%%  ---------------------------------------------------------------------------
%%  Muestra lo nuevo de la versión 2:
%%    · capas (fondo, frente y una capa propia) con elementos en posición absoluta
%%    · filas de columnas que terminan a la misma altura y llenan la página
%%    · cajas con 12 estilos, cuyo título se puede mostrar u ocultar
%%    · cajas de texto (T15–T19, T27, T29, T32, T35) con título opcional
%%  Compila con pdfLaTeX (Overleaf y latexmk repiten las pasadas solos).
%% ============================================================================
\documentclass{article}
\usepackage{postercientifico}

%% ---- Lienzo, letra, colores y estilo de las cajas ---------------------------
\posterconfig{
  tamano = A1, margen = 20mm,                % A1 vertical: 594 × 841 mm
  columnas = 3, separacion = 14mm,
  fuente = 27pt, familia = helvetica,
  idioma = spanish,                          % separación de palabras: spanish, english, portuguese…
  tema = ciruela,                            % sobrios: pizarra, salvia, niebla, ciruela, grafito;
                                             % vivos: azul, verde, burdeos, morado, gris, arena, oscuro
  cabecera = degradado, alineacion cabecera = izquierda,
  estilo bloque = caja,                      % barra, linea, caja, simple, borde, sombra, ...
  numerar bloques = true,
  mostrar titulos = true,                    % false oculta los títulos de todas las cajas
}

\titulo{Un póster por capas}
\subtitulo{Fondo, contenido y frente: cada capa se oculta, se atenúa y se ordena}
\autores{Nombre Apellido\textsuperscript{1}, Nombre Apellido\textsuperscript{2}}
\afiliaciones{\textsuperscript{1}Institución \quad \textsuperscript{2}Institución}
\logoderecho{img/logo1.png}

%% ---- Capas --------------------------------------------------------------------
%% Detrás del contenido (z < 0): una imagen muy tenue y dos formas de color
\colocarimagen[capa=fondo, x=0, y=0, ancla=north west, ancho=1, alto=1, ajuste=cubrir, opacidad=0.05]{img/prueba2.png}
\colocarforma[capa=fondo, forma=circulo, x=0.98, y=0.04, ancho=0.32, fondo=pcSecundario!14]
\colocarforma[capa=fondo, forma=rectangulo, x=0, y=1, ancla=south west, ancho=1, alto=0.012, fondo=pcAcento]

%% Una capa propia entre el contenido y el frente
\nuevacapa[z=5]{notas}
\colocartexto[capa=notas, x=0.632, y=0.262, ancla=north east, ancho=0.19, rotacion=3, fondo=FFF1A8, borde=E3C341,
  relleno=0.6em, alineacion=izquierda, escala=0.8, color=4A3B00]{\textbf{Nota adhesiva.} Está en la capa
  «notas» (z = 5): encima del contenido y debajo del frente.}

%% Delante de todo (z ≥ 0)
\colocartexto[capa=frente, x=0.5, y=0.55, rotacion=30, escala=9, opacidad=0.10, color=pcAcento, negrita]{BORRADOR}
%% Para quitar la marca de agua sin borrarla: \ocultarcapa{frente} oculta toda la capa;
%% o crea una capa solo para ella y ocúltala:  \nuevacapa[z=20, visible=false]{borrador}

\begin{document}
\cabecera

\begin{fila}[columnas=3, alto=llenar]
  \begin{columna}[distribucion=arriba]
    \begin{caja}[titulo={Introducción}]
      \txtEntradilla{Una frase que resume el mensaje del póster.}{\lipsum[1][1-7]}
    \end{caja}
    \begin{caja}[titulo={Objetivos}, estilo=linea]
      \txtNumerada{\item \lipsum[2][1-2] \item \lipsum[2][3-4] \item \lipsum[2][5-6]}
    \end{caja}
    \begin{caja}[titulo={Este título no se ve}, mostrar titulo=false, estilo=simple, estirar]
      \textbf{Caja sin título.} El título se guarda en el código (\texttt{mostrar titulo=false}) por si quieres
      recuperarlo. \lipsum[3][1-3]
      \imgPieLateral[0.5]{img/prueba4.png}{Una imagen con el pie al lado, centrado en vertical.}
      \lipsum[3][4-7]
    \end{caja}
  \end{columna}
  \begin{columna}[distribucion=arriba]
    \begin{caja}[titulo={Métodos}]
      \lipsum[4][1-3]
      \imgAncho{img/prueba1.png}{Figura principal de los métodos.}
      \txtResaltado[titulo={Idea clave}, mostrar titulo=true]{Un resaltado con su propio título. \lipsum[5][1]}
      \txtRecuadro{Un recuadro sin título (así es por defecto). \lipsum[5][2]}
    \end{caja}
    \begin{caja}[titulo={Resultados}, estilo=borde, estirar]
      \txtCifras{87\,\%}{precisión}{1\,234}{muestras}{0.92}{AUC}
      \imgDos{img/prueba6.png}{img/prueba3.png}{Dos paneles a la misma altura.}
      \txtTabla{Tabla 1. Resultados de relleno.}{%
        \begin{tabular}{lrrr}\toprule
          Grupo & Precisión & Recall & $n$\\\midrule
          Alfa & 0.91 & 0.88 & 134\\ Beta & 0.87 & 0.84 & 98\\ Gamma & 0.93 & 0.90 & 211\\\bottomrule
        \end{tabular}}
      \lipsum[6][1-4]
    \end{caja}
  \end{columna}
  \begin{columna}[distribucion=arriba]
    \begin{caja}[titulo={Discusión}, estilo=resaltado]
      \lipsum[7][1-6]
      \txtCaja[estilo=nota, titulo={Limitaciones}]{\lipsum[8][1-3]}
      \txtCaja[estilo=nota, titulo={Oculto}, mostrar titulo=false]{La misma caja, con el título oculto.}
      \txtLineaTiempo[titulo={Próximos pasos}, mostrar titulo=true]{
        \hito{2026}{Ampliar el conjunto de datos.}
        \hito{2027}{Validación externa.}
        \hito{2028}{Publicación del modelo.}}
    \end{caja}
    \begin{caja}[titulo={Conclusiones}, estilo=invertido, estirar]
      \txtVinetas{\item \lipsum[9][1-2] \item \lipsum[9][3-4] \item \lipsum[9][5-6]}
      \vfill
      \txtContacto[4.5em]{img/qr_prueba.png}{\textbf{Contacto}\par nombre.apellido\par @institucion.example}
    \end{caja}
  \end{columna}
\end{fila}

\piedeposter{Congreso de ejemplo \quad Ciudad, 2026}{}{Plantilla \texttt{postercientifico.sty} v1.0.0}
\end{document}
